% 2007 年新课标卷理科第 8 题：三视图（正视图、侧视图、俯视图）及尺寸。
%   正视图为直角边 20、20 的直角三角形；侧视图为底边 20、高 20 的等腰三角形；
%   俯视图为 20×20 正方形并画对角线，右侧标注 10 与 10。
% 来源 PDF：gaokaomath-src/普通高考/2007/2007新课标理(海南,宁夏).pdf
% SHA-256：c8ee1632ac3a7a28aab1150f047a2e1c74398ca2944701e9165237897ae724e2
% 原页：2007新课标理(海南,宁夏).pdf 第 1 页（A3 三栏）第 2 栏，书页码 393。
% 重绘依据：按原页三个视图的位置、形状与尺寸标注重绘，单位 cm。
\begin{tikzpicture}[x=0.115cm, y=0.095cm, semithick,
  every node/.style={inner sep=1.5pt, font=\small}]

  % ---------------- 正视图 ----------------
  \begin{scope}[shift={(0cm,0cm)}]
    \draw (0,0) -- (20,0) -- (20,20) -- cycle;
    % 尺寸 20（下）
    \draw[<->,>=stealth] (0,-4) -- (20,-4);
    \draw (0,-1.5) -- (0,-5.5);
    \draw (20,-1.5) -- (20,-5.5);
    \node[below] at (10,-4) {$20$};
    % 尺寸 20（右）
    \draw[<->,>=stealth] (24,0) -- (24,20);
    \draw (22,0) -- (26,0);
    \draw (22,20) -- (26,20);
    \node[right] at (24,10) {$20$};
    \node[below] at (10,-9) {正视图};
  \end{scope}

  % ---------------- 侧视图 ----------------
  \begin{scope}[shift={(3.6cm,0cm)}]
    \draw (0,0) -- (20,0) -- (10,20) -- cycle;
    \draw[<->,>=stealth] (0,-4) -- (20,-4);
    \draw (0,-1.5) -- (0,-5.5);
    \draw (20,-1.5) -- (20,-5.5);
    \node[below] at (10,-4) {$20$};
    \node[below] at (10,-9) {侧视图};
  \end{scope}

  % ---------------- 俯视图 ----------------
  \begin{scope}[shift={(0.9cm,-3.6cm)}]
    \draw (0,0) rectangle (20,20);
    \draw (0,0) -- (20,20);
    \draw (0,20) -- (20,0);
    % 尺寸 20（下）
    \draw[<->,>=stealth] (0,-4) -- (20,-4);
    \draw (0,-1.5) -- (0,-5.5);
    \draw (20,-1.5) -- (20,-5.5);
    \node[below] at (10,-4) {$20$};
    % 右侧 10 与 10
    \draw[<->,>=stealth] (24,0) -- (24,10);
    \draw[<->,>=stealth] (24,10) -- (24,20);
    \draw (22,0) -- (26,0);
    \draw (22,10) -- (26,10);
    \draw (22,20) -- (26,20);
    \node[right] at (24,5) {$10$};
    \node[right] at (24,15) {$10$};
    \node[below] at (10,-9) {俯视图};
  \end{scope}
\end{tikzpicture}
